\section*{الفصل \ref{ch:g6:shapes} --- المثلثات والرباعيات}

\begin{solution}{exo:g6:shapes:1}
اِتّبع الطريقة: \omterm{def:g6:lines:objects}{القطعة} $[AB]$ طولها $6$ سم، وقوس نصف قطره $5$ سم
مركزه $A$، وقوس نصف قطره $4$ سم مركزه $B$؛ ونقطة
تقاطعهما هي $C$.
\end{solution}

\begin{solution}{exo:g6:shapes:2}
القاعدة $[EF]$ طولها $3$ سم، ثم قوسان نصف قطرهما $5$ سم مركزهما $E$
و $F$؛ وتقاطعهما هو $D$. والزاويتان عند القاعدة $\widehat{DEF}$
و $\widehat{DFE}$ متساويتان (نحو $72.5^\circ$ لكل منهما) --- بتماثل
\omterm{def:g6:shapes:triangles}{المثلث المتساوي الساقين} (\cref{ex:g6:symmetry:proves}).
\end{solution}

\begin{solution}{exo:g6:shapes:3}
\omterm{def:g6:lines:objects}{قطعة} طولها $4$ سم، ثم قوسان نصف قطرهما $4$ سم مركزهما
طرفاها. وكل زاوية قياسها $60^\circ$.
\end{solution}

\begin{solution}{exo:g6:shapes:4}
(أ) \omterm{def:g6:shapes:triangles}{متساوي الأضلاع} (ومنه \omterm{def:g6:shapes:triangles}{متساوي الساقين} أيضًا). \quad
(ب) \omterm{def:g6:shapes:triangles}{متساوي الساقين} ($5 = 5$). \quad
(ج) \omterm{def:g6:shapes:triangles}{قائم الزاوية} \omterm{def:g6:shapes:triangles}{ومتساوي الساقين} (زاويتان متساويتان قياسهما $45^\circ$ تقابلان
ضلعين متساويين). \quad
(د) \omterm{def:g6:shapes:triangles}{متساوي الأضلاع} (ثلاث زوايا متساوية).
\end{solution}

\begin{solution}{exo:g6:shapes:5}
الوتر هو الضلع المقابل \omterm{def:g3:shapes:rightangle}{للزاوية القائمة}، وهو دائمًا أطول
الأضلاع. وفي المثلث $MNP$ \omterm{def:g6:shapes:triangles}{القائم الزاوية} في $N$، الوتر هو $[MP]$: فهو لا
يلمس أبدًا رأس \omterm{def:g3:shapes:rightangle}{الزاوية القائمة} $N$.
\end{solution}

\begin{solution}{exo:g6:shapes:6}
\emph{1. صواب}: فللمربع أربع \omterm{def:g3:shapes:rightangle}{زوايا قائمة}، وهو
تعريف المستطيل.

\emph{2. خطأ}: فالمستطيل $5 \times 3$ أضلاعه غير متساوية.

\emph{3. صواب}: فللمربع أربعة أضلاع متساوية، وهو تعريف
\omterm{def:g6:shapes:quadrilaterals}{المعيّن}.

\emph{4. صواب}: اُنظر \cref{ex:g6:shapes:recognize}.
\end{solution}

\begin{solution}{exo:g6:shapes:7}
إنشاء بالعموديات، كما في \cref{ex:g6:shapes:program}.
والقطران يقيسان نحو $5.8$ سم لكل منهما: فقطرا المستطيل
لهما الطول نفسه (\cref{prop:g6:shapes:diagonals}).
\end{solution}

\begin{solution}{exo:g6:shapes:8}
اُرسم \omterm{def:g6:lines:objects}{قطعة} طولها $6$ سم؛ ومنتصفها $O$؛ والعمودي في $O$؛
وضع عليه نقطتين على بُعد $2$ سم من $O$ في كل جهة. ووصل الأطراف
الأربعة يعطي \omterm{def:g6:shapes:quadrilaterals}{المعيّن} (\omterm{def:g5:solids:def}{فالرؤوس} الأربعة على المسافتين
$3$ سم و $2$ سم من $O$ على \omterm{def:g6:lines:objects}{مستقيمين} \omterm{def:g6:lines:perp}{متعامدين}، إذن الأضلاع الأربعة
متساوية بالتماثل).
\end{solution}

\begin{solution}{exo:g6:shapes:9}
برنامج صحيح:
1. اُرسم \omterm{def:g6:lines:objects}{القطعة} $[BC]$ حيث $BC = 4$ سم؛
2. اُرسم القوس الذي مركزه $B$ ونصف قطره $6$ سم فوق $[BC]$؛
3. اُرسم القوس الذي مركزه $C$ ونصف قطره $6$ سم فوق $[BC]$؛
4. سمِّ $A$ نقطة تقاطع القوسين؛
5. اُرسم \omterm{def:g6:lines:objects}{القطعتين} $[AB]$ و $[AC]$.
\end{solution}

\begin{solution}{exo:g6:shapes:10}
مهما كان موضع $M$ على الدائرة، تقيس الزاوية $\widehat{AMB}$
$90^\circ$: فالمثلث $AMB$ \omterm{def:g6:shapes:triangles}{قائم الزاوية} في $M$
كلما كانت $[AB]$ قطرًا. والبرهان ينتظر
\cref{ch:g8:cosine}.
\end{solution}

\begin{solution}{exo:g6:shapes:11}
قطران يتقاطعان في منتصفهما ولهما الطول نفسه:
تلك هي قائمة خاصّيات \emph{المستطيل}
(\cref{prop:g6:shapes:diagonals}). ولا شيء يفرض أن يكون القطران
\omterm{def:g6:lines:perp}{متعامدين}، فلا يلزم أن يكون معيّنًا ولا مربعًا --- والمستطيل
هو الجواب الصحيح (والمربع يحقّقها أيضًا، لكنه غير
مفروض).
\end{solution}

\begin{solution}{exo:g6:shapes:12}
مع $4 \times 4$ نقطًا توجد $3 \times 3$ خلايا في الشبكة.
والمربعات الموازية للشبكة: أضلاعها $1$ أو $2$ أو $3$ تعطي $9 + 4 + 1 = 14$
مربعًا. والمربعات المائلة موجودة أيضًا: فالصغيرة منها أضلاعها
أقطار خلايا $1 \times 1$ (مثل $(1,0)$ و $(2,1)$ و $(1,2)$ و
$(0,1)$) --- وعددها أربعة؛ والكبيرة أضلاعها
أقطار مستطيلات $2 \times 1$، مثل $(1,0)$ و $(3,1)$ و $(2,3)$ و
$(0,2)$ --- وعددها اثنان. وأضلاع مربع كهذا الأربعة متساوية
لأنها أقطار مستطيلات متطابقة. \omterm{def:g1:addition:def}{والمجموع}:
$14 + 4 + 2 = 20$ مربعًا. (وأيّ عدّ صحيح مع مثال مائل صحيح
ومع المربعات الموازية $14$ نجاح في هذا المستوى.)
\end{solution}

\begin{solution}{pb:g6:shapes:1}
\textbf{1.} الأضلاع الأربعة كلها بالطول نفسه (نحو
$4.2$ سم) والزوايا الأربع قائمة: فالشكل
\emph{مربع}. (فتساوي القطرين يعطي \omterm{def:g3:shapes:rightangle}{الزوايا القائمة}، كما في
المستطيل؛ وتعامدهما يعطي تساوي الأضلاع، كما في
\omterm{def:g6:shapes:quadrilaterals}{المعيّن}؛ والأمران معًا يعطيان المربع، وهو
\cref{prop:g6:shapes:diagonals} مقروءًا بالمقلوب.)

\textbf{2.} \emph{\omterm{def:g6:shapes:quadrilaterals}{معيّن}}: أربعة أضلاع متساوية (نحو $4.7$ سم)،
لكن لا \omterm{def:g3:shapes:rightangle}{زاوية قائمة} --- وهو بالضبط إنشاء
\cref{exo:g6:shapes:8}.

\textbf{3.} تخرج الزوايا الأربع قائمة لكن الأضلاع تأتي
بطولين مختلفين: \emph{مستطيل} --- وهي
حالة \cref{exo:g6:shapes:11}.

\textbf{4.} لا يصمد أيّ ادّعاء بتساوي الأطوال أمام القياس ولا
تظهر \omterm{def:g3:shapes:rightangle}{زاوية قائمة}؛ لكن الأضلاع المتقابلة متساوية مثنى مثنى
وتبدو متوازية. وهو «المستطيل المائل» العامّ الذي
ينصّف قطراه كل منهما الآخر لا غير --- أي
\emph{متوازي الأضلاع}، ويُسمّى ويُدرس في \cref{ch:g7:central}.

\textbf{5.} قطران يتقاطعان في منتصفهما المشترك:
\begin{center}
\renewcommand{\arraystretch}{1.2}
\begin{tabular}{l|l}
القطران\dots & الشكل \\
\hline
متساويان ومتعامدان & مربع \\
متعامدان فقط & \omterm{def:g6:shapes:quadrilaterals}{معيّن} \\
متساويان فقط & مستطيل \\
لا هذا ولا ذاك & الشكل الذي بلا اسم (اُنظر \cref{ch:g7:central}) \\
\end{tabular}
\end{center}

\textbf{6.} لكل مثلث ركن ضلع قائمة على كل
\omterm{def:g4:geometry:circle}{قطر}: نصف $8$ سم ونصف $5$ سم، أي ضلعان قائمتان $4$ سم
و $2.5$ سم، بينهما \omterm{def:g3:shapes:rightangle}{زاوية قائمة} --- والمعطيات نفسها
للأربعة جميعًا. والمثلثات المبنية من الضلعين نفسيهما ومن
\omterm{def:g3:shapes:rightangle}{الزاوية القائمة} بينهما نسخ مطابقة بعضها لبعض
(اِبنِها: فضبطات الفرجار متطابقة). ومنه تكون أوتارها
الأربعة --- وهي أضلاع الشكل --- متساوية:
فالشكل \omterm{def:g6:shapes:quadrilaterals}{معيّن}.

\textbf{7.} اِفتح الفرجار على طول ضلع واحد، ثم
اِمشِ به حول الشكل: فإن وقعت سنّه وقلمه بالضبط
على كل زوجين من الأركان المتجاورة، فالأضلاع الأربعة كلها تساوي تلك
الفتحة. ولا حاجة إلى مسطرة.

\textbf{8.} البرنامج:
\begin{enumerate}
  \item اُرسم \omterm{def:g6:lines:objects}{القطعة} $[AC]$ حيث $AC = 6$ سم وضع
        منتصفها $O$ (قِس $3$ سم من $A$).
  \item اُرسم العمودي على $(AC)$ في $O$.
  \item وضع عليه $B$ و $D$ على بُعد $3$ سم من $O$، واحدة في كل
        جهة.
  \item صِل $A$ و $B$ و $C$ و $D$.
\end{enumerate}
القطران متساويان ($6$ سم)، ومتعامدان، وينصّف كل منهما الآخر:
فحسب القاموس، $ABCD$ مربع.

\textbf{9.} يعطي القياس نحو $4.2$ سم: أي \emph{أطول} من
$3$ سم --- فالتخمين الشائع «\omterm{def:g3:shapes:shapes}{نصف القطر}» خطأ. وأما
النصف القابل للبرهنة: فللانتقال من الركن $A$ إلى الركن التالي $B$، يمكن
المشي على الضلع $[AB]$ مستقيمًا، أو أخذ الالتفاف
$A \to O \to B$ عبر المركز، وثمنه $3 + 3 = 6$ سم؛ والضلع
\omterm{def:g6:lines:objects}{المستقيم} أقصر، إذن الضلع أقلّ من $6$ سم.
وأما أنه \emph{أكثر} من $3$ سم فهو ما تبيّنه المسطرة؛
وسيبرهن \cref{ch:g8:pythagoras} عليه ويعطي القيمة المضبوطة
($3$ سم مضروبًا في الجذر التربيعي للعدد $2$، أي نحو $4.24$ سم).

\textbf{10.} مثلثات الأركان الأربعة عند $H$ تأتي الآن في
زوجين \emph{متطابقين}: اثنان ضلعا قائمتهما $3$ سم و $2.5$ سم
(وتراهما $[AB]$ و $[AD]$)، واثنان ضلعا قائمتهما $5$ سم و
$2.5$ سم (وتراهما $[CB]$ و $[CD]$). إذن $AB = AD$ و
$CB = CD$: فللطائرة الورقية زوجان من الأضلاع \emph{المتجاورة}
المتساوية (\omterm{def:g6:shapes:quadrilaterals}{والمعيّن} هو \omterm{pb:g6:shapes:1}{الطائرة الورقية} الخاصّة التي تتطابق فيها الأربعة كلها).

\textbf{11.} \omterm{def:g4:geometry:polygon}{الرباعي}: $2$ قطران. \omterm{def:g4:geometry:polygon}{والخماسي}: $5$.
\omterm{def:g4:geometry:polygon}{والسداسي}: $9$.

\textbf{12.} من رأس واحد \omterm{def:g4:geometry:polygon}{للسداسي}، تذهب الأقطار إلى
كل رأس عدا نفسه وجاريه: أي $6 - 3 = 3$
أقطار. \omterm{def:g5:solids:def}{والرؤوس} الستة تطلق $6 \times 3 = 18$ طرف \omterm{def:g4:geometry:circle}{قطر}
--- لكن كل \omterm{def:g4:geometry:circle}{قطر} عُدّ عند \emph{طرفيه} معًا،
مرة من كل طرف، تمامًا كمصافحة يعدّها
المتصافحان معًا (\cref{pb:g6:wholes:1}). والعدّ الصحيح:
$18 \div 2 = 9$. وهو يوافق الرسم.

\textbf{13.} عشرة \omterm{def:g5:solids:def}{رؤوس}: $10 - 3 = 7$ من كل رأس، إذن
$10 \times 7 \div 2 = 35$ قطرًا. وعشرون رأسًا:
$20 \times 17 \div 2 = 170$ قطرًا.

\textbf{14.} جرّب المرشّحين: $10$ \omterm{def:g5:solids:def}{رؤوس} تعطي $35$ (السؤال
13)؛ و $11$ رأسًا تعطي $11 \times 8 \div 2 = 44$. وُجد: فللمضلع
$11$ رأسًا.

\textbf{15.} ليس للمثلث أقطار البتة: فغير المجاورين
لكل رأس هم\,\dots\ لا أحد، لأن الرأسين الآخرين
جاراه كلاهما. والطريقة توافق ذلك:
$3 - 3 = 0$ من كل رأس، \omterm{def:g1:addition:def}{والمجموع} $0$. فالمصافحات
والأقطار هما العدّ نفسه في ثوب مختلف: أزواج تُختار
من مجموعة، كل زوج يُعدّ مرة --- والرياضيات تعيد استعمال
أفكارها الجيدة.
\end{solution}
